% Part: first-order-logic
% Chapter: introduction
% Section: first-order-logic

\documentclass[../../../include/open-logic-section]{subfiles}

\begin{document}

\olfileid[hi]{fol}{int}{fol}

\olsection{प्रथम-क्रम तर्क}

औपचारिक तर्कशास्त्र से अपने पहले परिचय के कारण आप संभवतः प्रथम-क्रम तर्क से
परिचित हैं।\footnote{वास्तव में हम कमोबेश यही मानकर चलते हैं!{} यदि आप
परिचित नहीं हैं, तो \emph{forall x} \citep{Magnus2021} जैसी कोई अधिक
आरंभिक पाठ्यपुस्तक दोहरा सकते हैं।} आप इसे ``परिमाणक तर्क'' अथवा ``विधेय
तर्क'' के नाम से जानते होंगे। सबसे पहले, प्रथम-क्रम तर्क एक औपचारिक भाषा है।
इसका अर्थ है कि उसकी एक निश्चित शब्दावली है और उसके व्यंजक उस शब्दावली से
बनी शृंखलाएँ हैं। किंतु प्रत्येक शृंखला अनुमत नहीं होती। अनुमत व्यंजकों की
कई श्रेणियाँ हैं: पदों, !!{formula}s और !!{sentence}s की। हमारी मुख्य रुचि
प्रथम-क्रम तर्क के !!{sentence}s में है: वे हमें सामान्य भाषा के वाक्यों का
औपचारिक अनुरूप देते हैं और उनके बारे में हम वे प्रश्न पूछ सकते हैं जिनमें
तर्कशास्त्रियों की सामान्यतः रुचि होती है। उदाहरणार्थ:
\begin{itemize}
    \item क्या $!B$,~$!A$ से तार्किक रूप से निष्पन्न होता है?
    \item क्या $!A$ तार्किक रूप से सत्य, तार्किक रूप से असत्य, अथवा
आकस्मिक है?
    \item क्या $!A$ और $!B$ समतुल्य हैं?
\end{itemize}

ये प्रश्न मुख्यतः प्रथम-क्रम तर्क के !!{sentence}s के ``अर्थ'' से जुड़े
हैं। उदाहरणार्थ, कोई दार्शनिक इस प्रश्न का कि क्या $!B$,~$!A$ से तार्किक
रूप से निष्पन्न होता है, इस प्रकार विश्लेषण करेगा: क्या कोई ऐसी स्थिति है
जिसमें $!A$ सत्य किंतु~$!B$ असत्य है ($!B$,~$!A$ से निष्पन्न नहीं होता), या
फिर $!A$ को सत्य बनाने वाली प्रत्येक स्थिति~$!B$ को भी सत्य बनाती है
($!B$,~$!A$ से निष्पन्न होता है)? किंतु अभी हमें यह नहीं बताया गया कि
``स्थिति'' क्या है---यह \emph{अर्थविज्ञान} का काम है। प्रथम-क्रम तर्क का
अर्थविज्ञान ``स्थिति'' की दार्शनिक सहज धारणा का गणितीय रूप से सटीक प्रतिरूप
देता है; साथ ही---और यह महत्त्वपूर्ण है---यह भी सटीक करता है कि किसी स्थिति
  में !!a{sentence}~$!A$ का \emph{सत्य होना} क्या है। जिस गणितीय रूप से सटीक
  प्रतिरूप का हम विकास करेंगे, उसे !!a{structure} कहते हैं। ``में सत्य'' को
  सटीक करने वाला संबंध \emph{संतुष्टि} संबंध कहलाता है। अतः हम ``$!A$,
  ~$\Struct{M}$ में संतुष्ट है'' (चिह्नों में: $\Sat{M}{!A}$) को
  !!{sentence}s के प्रत्येक सदस्य~$!A$ के लिए परिभाषित करेंगे; इसमें
  !!{structure}s~$\Struct{M}$ संभावित व्याख्याएँ होंगी।
  इसके बाद ``से निष्पन्न
होता है'' अथवा ``तार्किक रूप से सत्य है'' जैसे अन्य अर्थवैज्ञानिक पदों की भी
सटीक परिभाषाएँ दी जा सकेंगी। इन परिभाषाओं से हम गणितीय सटीकता के साथ यह
निर्धारित कर पाएँगे कि, उदाहरणार्थ, क्या $\lforall[x][(!A(x) \lif !B(x)),
\lexists[x][!A(x)] \Entails \lexists[x][!B(x)]]$। उत्तर निस्संदेह ``हाँ''
होगा। यदि आपने सामान्य भाषा के वाक्यों को प्रथम-क्रम तर्क में प्रतीकित करना
सीखा है, तो आप इसे, उदाहरणार्थ, ``सभी चींटियाँ कीट हैं; कुछ चींटियाँ हैं; अतः
कीट हैं'' के प्रतीकीकरण के रूप में पहचानेंगे। यह स्पष्टतः वैध तर्क है, इसलिए
हमारी औपचारिक भाषा में ``से निष्पन्न होता है'' का गणितीय प्रतिरूप भी यही
उत्तर देना चाहिए।

औपचारिक तर्कशास्त्र से अपने पहले परिचय का एक और विषय आपको याद होगा:
\emph{!!{derivation}s}। यदि आपने औपचारिक तर्कशास्त्र का आरंभिक पाठ्यक्रम
लिया है, तो आपके शिक्षक ने ऐसे उदाहरण खोजने का अभ्यास कराया होगा जो
!!{derivation}s हों, संभवतः ऐसी !!a{derivation} भी जो दिखाती है कि ऊपर का निष्पत्ति-संबंध सत्य
है। !!{derivation}s देने की कई अलग विधियाँ हैं: आपने ``प्राकृतिक निष्पत्ति''
अथवा ``सत्यता-वृक्ष'' नामक कोई विधि अपनाई होगी, किंतु और भी बहुत-सी हैं।
!!{derivation} प्रणालियों का उद्देश्य ऐसे साधन देना है जिनसे तर्कशास्त्रियों
के ऊपर के प्रश्नों का उत्तर दिया जा सके। उदाहरणार्थ, कोई प्राकृतिक निष्पत्ति
!!{derivation}, जिसमें $\lforall[x][(!A(x) \lif !B(x))$ और
$\lexists[x][!A(x)]$ आधार-वाक्य तथा $\lexists[x][!B(x)]]$ निष्कर्ष (अंतिम
पंक्ति) हैं, यह \emph{सत्यापित} करती है कि $\lexists[x][!B(x)]$, तार्किक रूप
से $\lforall[x][(!A(x) \lif !B(x))]$ और $\lexists[x][!A(x)]$ से निष्पन्न
होता है।

पर ऐसा क्यों है? प्रथम दृष्टि में !!{derivation} प्रणालियों का अर्थविज्ञान
से कोई संबंध नहीं दिखता: औपचारिक !!{derivation} देने में केवल चिह्नों को
नियमों से नियंत्रित कुछ ढंगों में व्यवस्थित किया जाता है; उसमें ``स्थितियों''
अथवा ``में सत्य'' का कोई उल्लेख नहीं होता। !!{derivation} प्रणालियों और
अर्थविज्ञान के बीच संबंध को अधितार्किक जाँच से स्थापित करना पड़ता है।
उदाहरणार्थ, हमें ऐसा गणितीय प्रमाण चाहिए कि $\lexists[x][!B(x)]$ की औपचारिक
!!{derivation}, आधार-वाक्यों $\lforall[x][(!A(x) \lif !B(x))]$ और
$\lexists[x][!A(x)]$ से, तभी और केवल तभी संभव है जब $\lforall[x][(!A(x)
\lif !B(x))$ और $\lexists[x][!A(x)]$ मिलकर
$\lexists[x][!B(x)]]$ को निष्पन्न करते हों। किंतु ऐसा करने से पहले वाक्यविन्यास
और अर्थविज्ञान की परिभाषाओं को ठीक करने के लिए बहुत-सा सावधानीपूर्ण कार्य
करना होगा।

\end{document}
